\documentclass[UTF8,fontset=fandol,11pt,a4paper]{ctexart}
\usepackage[top=16mm,bottom=16mm,left=20mm,right=20mm]{geometry}
\usepackage{amsmath,booktabs,tabularx,array,graphicx,caption,float,setspace,hyperref,needspace}
\hypersetup{hidelinks,pdftitle={因子流动性实验：2024年连续账户与资金阶梯},pdfauthor={factor-liquidity}}
\setstretch{1.04}\setlength{\parskip}{0.15em}\setlength{\parindent}{2em}
\setlength{\textfloatsep}{7pt}\setlength{\intextsep}{6pt}
\captionsetup{font=small,labelfont=bf,labelsep=quad,skip=4pt}
\newcommand{\lead}[1]{\par\medskip\Needspace{3\baselineskip}\noindent\textbf{#1}\par\smallskip}
\newcommand{\panel}[2]{\begin{figure}[H]\centering\includegraphics[width=0.98\linewidth]{figures/#1.pdf}\caption{#2}\label{fig:#1}\end{figure}}
\clubpenalty=10000\widowpenalty=10000\raggedbottom
\begin{document}
\noindent{\large\bfseries 因子流动性实验}\hfill{\small 2026 年 10 月 26 日}
\lead{实验结论}
本次已完成原 G1 因子在2024年1月2日至8月30日的连续账户回测，覆盖162个交易日、8次新目标和25组资金与成交情景，全部通过账户核账。以1亿元、10\%参与率为例，全期收益为 $-10.7623\%$，同资金参考账户为 $-9.4077\%$，两者相差 $-1.3546$ 个百分点；新目标日可定量计划的当日完成率为99.6426\%。费用和税费均为零。

这些结果表明，大部分已经定量的计划金额能在当天完成，但原组合的覆盖、剩余小额需求和实际资金投入仍需一起看。首期956个原目标中，875个可以定量，日末实际持有873只证券。100万元账户则只建仓327只，首日现金比例为72.70\%；资金规模改变的既有成交约束，也有分散持仓能否满足最小申报数量的问题。

\lead{因子、数据和实验设置}
沿用原 N20、sample3、\texttt{H20\_equal\_G1}，取原处理后低因子 G1，保持原权重和 H20 相位。正式因子的原始公式是 $f_t=C_t/C_{t-20}-1$，其中 $C$ 为原批次复权收盘价。之后经过缺失处理、去极值、行业和市值中性化及标准化，才分成五组；G1是处理后因子的低分组，不能直接理解为过去跌幅最大的股票。$N=20$是回看期，$H=20$是调仓间隔。

正式组合使用信号日当时的历史100001股票池。目标生成沿用原来的流动性筛选：最近20个正成交额观测的均值严格大于500万元。这是来源策略已有的条件，不是本次新增的容量合格线。调仓直接采用上一交易日已经确定的筛选结果，不再额外延迟一个交易日；股票池退出时目标归零，退出权重不重新分配。

行情使用公司TDengine的\texttt{level2.stb\_depthmarket}，按未复权的元／股或元／份、股或份解释。官方日历、涨跌停价、停牌和现有公司行动资料用于共同约束。日线只用于已到日终的估值和库内交叉核对，不作为早晨已知价格。此前合成案例和较新市场小样本用于检查执行程序，本报告的数值均来自这次2024年真实行情账户。

五档初始资金为100、300、1000、3000和10000万元，每档分别运行参考账户、L2不限参与率及5\%、10\%、20\%参与率，共25组。每个账户都在1月2日一次投入现金，从零持仓、零历史余单开始；1月2日至10日持现金，1月11日首次建仓。随后沿原相位在2月8日、3月15日、4月16日、5月17日、6月17日、7月15日和8月12日更新目标。

\lead{账户和成交怎样模拟}
每天09:35先处理已经生效的权益，用当时已知价格估值并冻结该账户净资产，再一次生成全部证券当天的买卖计划。新计划的参考价必须严格早于09:35，且年龄不超过30秒；已有持仓缺少新报价时，按后文约定沿用估值。已经完成的目标不因每日价格变化重新平衡；未完成方向次日按实际持仓重新定量，新目标或退出等明确事件才重新判断方向。执行时段为[09:35,11:30)和[13:00,14:57)，新子单必须经过正延迟，在后续合格快照成交。

参考账户以事先确定的参考价作为反事实成交价，不限盘口数量和参与率，但保留现金、数量、T+1、停牌、权益、延迟和共同数据资格约束。L2账户按可见对手盘价格和数量模拟成交。区间均价（VWAP）是该区间成交额除以成交量，成交量为零时不计算；订单成交均价则只用真正模拟成交的价格与数量加权。例如，卖一10元有200股、卖二10.01元有100股，买入这300股花费3001元，均价约10.0033元，不能直接用市场区间均价替代。

一个账户只有一份现金余额。假设余额10万元，已有买单预占6万元，其他买单只能使用余下4万元；撤单只释放预占，不增加现金余额。实际卖出释放现金后才能用于买入，未到账红利不能使用。同一市场时点先处理可执行卖单，再处理已有合法预占的买单，随后生成新子单；新买单不能无延迟吃掉用于决策的同一张快照。

同证券的多笔任务还共享一份参与率额度。例如，10\%情景在需求激活后已观察到1万股有效成交量，最多获得1000股累计额度；第一笔用了600股，其他任务合计最多再用400股。只有已经结束、被观察且质量合格的区间才授额度，不能领取未来量。超过10秒的无法解释间隔、量额异常或累计重置会使相关连续段暂停并重建基准，不跨断点补回额度。不同资金与情景只共享只读行情，各自保存现金、持仓、预占和额度。

\lead{指标怎样计算}
净资产等于真实现金、实际股份市值、现金应收及尚未记入实际股份的应收股份市值之和；本研究不融资。应收到账只是科目移动，不再次增加净资产。预占仍属于现金余额，只限制可用现金，不重复扣减净资产。净值为日末净资产除以唯一初始资金，最大回撤包含共同现金起点。
\[
\text{计划完成率}=\frac{\sum(\text{已成交数量}\times\text{本父单参考价})}
{\sum(\text{计划数量}\times\text{本父单参考价})}.
\]
买卖金额均按正数汇总，现金不足不缩小分母。主图只统计各次新目标日09:35固定的计划，次日剩余需求不再次算作新计划。缺参考目标数量未定，另列原权重和理论金额。日末持仓差距是已定量目标与实际经济持股之差的绝对值，按各参考价折算后除以冻结净资产；经济持股包括已生效的应收股份。

滑点只计算模拟成交部分：买得更贵或卖得更便宜记为正，按成交数量乘参考价加权，以基点报告。它包含等待期间的市场变化，不能全部称为自身市场冲击。完成耗时只统计新目标当天已完成的任务，从09:35算到最后一笔成交，交易分钟不含午休和隔夜；未完成任务同时报告数量及剩余金额，不补造完成时间。

\lead{1亿元账户的连续结果}
47日和67日是同一条162日连续路径的中间观察，没有重置现金或拼接独立窗口。表中收益差均为L2账户减同资金参考账户。
\begin{table}[H]\centering\small
\begin{tabular}{llrrrr}\toprule
窗口截至 & 情景 & 收益（\%） & 收益差（百分点） & 当日完成（\%） & 最大回撤（\%）\\\midrule
03-14 & 参考账户 & 4.2990 & 0.0000 & 99.8610 & 26.2769 \\
03-14 & L2 10\% & 3.8689 & -0.4301 & 99.6459 & 26.3348 \\
04-15 & 参考账户 & 0.9929 & 0.0000 & 99.8970 & 26.2769 \\
04-15 & L2 10\% & 0.3529 & -0.6401 & 99.7286 & 26.3348 \\
08-30 & 参考账户 & -9.4077 & 0.0000 & 99.8239 & 26.2769 \\
08-30 & L2 10\% & -10.7623 & -1.3546 & 99.6426 & 26.3348 \\
\bottomrule\end{tabular}
\caption{1亿元、参考与10\%情景的同口径连续窗口。完成率只针对已定量的新目标日计划。}
\end{table}
\panel{01_research_nav}{1亿元五种情景的连续净值与回撤。每条路径独立承接现金、持仓和权益，所有官方交易日均保留。}

10\%账户在2月5日出现全期最深回撤26.3348\%，当日净值为0.7459；3月14日累计收益回升至3.8689\%，到8月30日净值为0.8924。参考账户同期末净值为0.9059。因此，全期既包含持有期间的市场变化，也包含两种执行路径之间的累计差异；后者对应期末资产相差135.46万元。

\lead{八次调仓的完成进度}
1月11日首次建仓，10\%账户在30分钟、60分钟和当日结束时，分别完成99.2566\%、99.4425\%和99.4425\%的当次计划金额。4月16日的当日完成率为99.3479\%，剩余计划参考金额103.92万元，已定量持仓差距占冻结净资产1.0532\%。这两个比例分母不同：前者比较买卖计划金额，后者比较账户净资产。
\panel{02_new_target_progress}{1亿元各次新目标的30分钟、60分钟、当日进度，以及已定量目标的日末持仓差距。次日重试不回写当天进度。}

汇总八次新目标，10\%账户可定量计划的当日完成率为99.6426\%，其中1668笔当天未完全成交，剩余参考金额合计399.42万元。按金额看大部分需求已完成，仍可能留下较多小额任务；这不等于它们到窗口末都未完成。8月30日仍有78个正权重目标的执行方向尚未结束，其中也包括数据或数量限制。
\begin{table}[H]\centering\small
\begin{tabular}{lrrrrr}\toprule
情景 & 当天完成数 & 平均分钟 & 90分位分钟 & 当天未完数 & 未完金额（万元）\\\midrule
参考账户 & 11,998 & 0.14 & 0.17 & 121 & 197.98 \\
L2 10\% & 10,443 & 4.31 & 10.45 & 1,668 & 399.42 \\
\bottomrule\end{tabular}
\caption{仅看新目标当天的固定任务。耗时只统计其中已完成者；未完金额按各父参考价折算，与收益损失不是同一指标。}
\end{table}
10\%账户全部模拟成交的加权滑点为11.72个基点。参考路径与L2路径的实际持仓和净资产逐渐不同，因此后续计划数量也可不同；这是相同目标策略的连续账户对照，不是全期逐笔相同订单的价格比较。

\lead{资金规模怎样改变实际组合}
图\ref{fig:05_capital_staircase}给出全部25组结果。初始资金较小时，分到每只证券的金额可能不足合法申报数量，账户会保留大量现金；资金增加后，持仓覆盖和投入比例提高，盘口与参与率约束的作用也随之变化。
\begin{table}[H]\centering\small
\begin{tabular}{rrrrrr}\toprule
初始资金（万元） & 参考账户 & 不限参与率 & L2 5\% & L2 10\% & L2 20\%\\\midrule
100 & 1.1594 & 0.3391 & 0.2554 & 0.2651 & 0.2776 \\
300 & -2.9042 & -3.8293 & -3.8269 & -3.8818 & -3.8882 \\
1,000 & -7.8525 & -8.9550 & -8.6586 & -8.7643 & -8.8384 \\
3,000 & -9.0908 & -10.3668 & -10.3192 & -10.3587 & -10.3054 \\
10,000 & -9.4077 & -11.2410 & -10.6392 & -10.7623 & -10.8213 \\
\bottomrule\end{tabular}
\caption{全部25组的全期收益，单位为\%。费用、税费均为零。}
\end{table}
\panel{05_capital_staircase}{五档资金下的全期收益和新目标日当日完成率。每档都有自己的参考账户，不能用一个资金档的参考收益替代其他档。}

以10\%情景为例，100万元账户首期每个原目标理论上只分到约1046元，取整后453个目标有正数量，实际建仓327只，现金占净资产72.70\%；1亿元对应875个正数量目标、873只持仓和9.48\%现金。正数量仍可能低于板块最小申报要求，因此取整后目标数与实际建仓数也不必相等。
\begin{table}[H]\centering\small
\begin{tabular}{rrrrr}\toprule
初始资金（万元） & 原正目标数 & 取整后正目标数 & 实际建仓数 & 首日现金比例（\%）\\\midrule
100 & 956 & 453 & 327 & 72.70 \\
300 & 956 & 774 & 659 & 42.34 \\
1,000 & 956 & 870 & 841 & 18.95 \\
3,000 & 956 & 875 & 872 & 11.69 \\
10,000 & 956 & 875 & 873 & 9.48 \\
\bottomrule\end{tabular}
\caption{五档资金、10\%情景在1月11日的覆盖与投入。原权重保留，没有把受阻证券的权重分给其他证券。}
\end{table}

100万元、10\%账户全期收益0.2651\%，当日完成率85.3915\%；1亿元对应$-10.7623\%$和99.6426\%。两者同时包含不同的实际投入比例与持仓覆盖，不能只据绝对收益高低判断资金容量。参与率提高也不保证收益单调改善：1亿元不限参与率的完成率为99.8202\%，收益为$-11.2410\%$，而10\%情景分别为99.6426\%和$-10.7623\%$。成交时点和实际持仓路径会改变收益，完成率与收益需要分别看。

\lead{现金、缺数据和未成交原因}
缺少合格早晨参考的证券当天不编造数量，保留目标并次日重试；有持仓后缺L2，只暂停相关执行。日线用于日终估值，日线也缺时沿用最近已知有效价并记录暴露。现金比例由实际成交、数量取整和持仓变化形成，不强制等于缺数权重对应的金额。
\panel{03_data_cash_rights}{1亿元建仓后的现金比例，以及10\%情景中的数据、数量、陈旧估值和应收权益暴露。各限制按自身定义统计，不直接相加为组合缺口。}

10\%账户首日现金比例为9.4789\%，期末为6.8312\%；每日缺参考目标权重最高为0.4162\%。沿用估值的账面金额占净资产最高为1.6593\%，出现在4月30日。日末现金应收最高103.80元，股份应收日末余额始终为零；应收不作为已到账现金或可卖股份使用。

新目标当天剩余计划按主原因分类：行情或量源检查238.27万元、数量单位159.85万元、参与率额度1.30万元。这里的主原因说明程序最后在哪项条件下保留了余量，同一任务之前仍可能受到其他限制；不能把行情检查的限制全部解释成市场本身没有成交能力。
\panel{04_unfilled_reasons}{1亿元、10\%情景的新目标日未完成计划，每笔按固定检查顺序归入一个主原因；不重复计入多类，也不累加次日重试。}

\lead{结果适用范围}
本批次对关键权益资料不足以独立确定股份分配等经济事实的215只证券采用共同限制清单，所有资金和情景从相同现金起点不交易它们，保留原权重、不重新分配。1亿元、10\%路径的原目标中，该清单限制权重最高为9.0531\%。这是按全窗口资料统一确定的数据条件，不是当时可用的选股筛选，也不表示这些证券当时不能交易。

本轮沿用已保存的行情批次。10月26日15:32重新查询时，旧清单581个整日零L2证券日都已出现记录；这些新记录未自动替换本轮缓存，因此本轮中的689009缺参考仍按原研究规则暂停。源库已经补齐与本轮实际使用哪些数据，是两个不同事实。历史回放数据可修订，本结果也不声称还原了当时数据库的每个历史版本。

公司行动按现有事件资料记账，缺全市场完整覆盖声明记为警告。只有日期而无可靠时刻时，登记在日末锁定资格，明确的除息／权益生效日在09:35前确认经济权益，付款、股份到账和可卖日期在已知日期结束后更新，下一交易日使用。到账日未知但权益价值明确时保留应收；股份已到账但可卖日期未知时记入实际持仓，暂不计入可卖数量。陈旧价跨越明确的简单权益只做一次带标记的理论除权估值，不把它当成交价，也不叠加复权收益。

本实验回答原G1目标在共同数据条件、既定现金与权益规则下怎样执行。合格快照刷新可见深度、以正延迟后的后续快照成交都是模拟约定，不还原真实队列，也不测自身市场冲击。2024年参与过原策略参数选择，本轮不称为独立样本外；未计算同期100001价格指数超额，不以其他指数替代。最低收益要求尚未确定，本报告不据此给出业务容量上限。
\end{document}
